\documentclass[20pt,margin=0in,innermargin=0in]{tikzposter}
\geometry{paperwidth=36in,paperheight=24in}
\usepackage[utf8]{inputenc}
\usepackage{amsmath}
\usepackage{amsfonts,sourceserifpro}
\usepackage{amsthm}
\usepackage{amssymb}
\usepackage{mathrsfs,blindtext,nth,microtype}
\usepackage{graphicx}
\usepackage{adjustbox}
\usepackage{enumitem}
\usepackage[backend=bibtex,style=numeric]{biblatex}
\usepackage{stanford-theme}

\usepackage{mwe} % for placeholder images
\bibliography{references}


% set theme parameters
\tikzposterlatexaffectionproofoff
\usetheme{StanfordTheme}
\usecolorstyle{StanfordStyle}


\title{Poster Template}
\author{person name\\\texttt{personname@stanford.edu}}
\institute{Department of Computer Science, Stanford University} % you probably don't need this at a Stanford event? Whatever lol
\titlegraphic{\includegraphics[height=0.125\textheight]{stanford}}


\begin{document}

\maketitle[titletotopverticalspace=1in]
\centering

\begin{columns}
    \column{0.3}
    \block{\nth{1} Column}{

    You can have \textbf{however many} columns as you want, and
    you can make them \textbf{whatever widths} you want.
    You can even have columns within columns.

    \vspace{1em}

    You should take this document as a \textbf{starting point}, if you take it at all.
    It's no good if everyone copies this layout, pastes in their text, and goes off to the printers.
    It's worth thinking for a minute or two about if you want the poster in thirds, or half-half, or what.
    
    \vspace{1em}

    To change the width of a column, change the value for \texttt{\textbackslash column\{0.3\}} (30\% of the width of the page) to something else (like \texttt{\textbackslash column\{0.4\}} for 40\% or \texttt{\textbackslash column\{0.62312\}} for 62.312\%). Just try to get the columns to add up to 1.0.


    }
    
    \block{\nth{2} Column}{
    \blindtext[1]

    \vspace{1em}

    {\scshape Don't get too close to margins.

    No, seriously. There's nothing but trouble here.}

    }

    \column{0.45}
    \block{\nth{3} Column}{
        You can have some formulas if you want:

        \begin{align*} q^{-3} & \le \frac{\overline{\sqrt{2}-\emptyset}}{\tilde{\omega} \left( e, \dots, \frac{1}{P ( A )} \right)} \wedge p \left( \bar{K}^{-5}, \tilde{m} \right) \\ & = \max_{B \to \emptyset}  1 \pm \dots \cup \pi \left(-q ( d ), \dots, \mathscr{{C}}'' \right)  \\ & \le \left\{ 1^{-7} \colon \cosh^{-1} \left(-\kappa \right) \le \max \int_{\hat{M}} \tanh \left( C^{5} \right) \,d \theta \right\} \\ & \le \prod  \cosh^{-1} \left( \pi^{-8} \right) + \dots \vee \omega \left(-\pi, \infty \sqrt{2} \right) \end{align*}

        You can also have a figure:
        
        \vspace{1em}
        \begin{tikzfigure}[Here's where captions go. Captions are good.]
            \includegraphics[width=0.8\linewidth]{example-image}
        \end{tikzfigure}
        % \vspace{1em}
    }




    \column{0.25}
    \block{\nth{4} Column}{
    I think you're getting the idea.
        
    \vspace{6em}
    }
    

    \block{Acknowledgments}{
    Be sure to thank people who made the work you did possible --- specifically, funding sources, people and organizations who provided data, etc\dots

    \vspace{2em}

    There's no need to acknowledge your collaborators (the people you list at the top of the poster), because their contribution is being acknowledged in the form of co-authorship.

    \vspace{2em}

    \textit{Fun fact: author order is ``a whole thing''; you should check with your mentor about the norms in whatever field you're publishing to.}
    }
    
    % \block{References}{
    %     \vspace{-1em}
    %     \begin{footnotesize}
    %     \printbibliography[heading=none]
    %     You probably won't need to list all the references, but if you do, you can uncomment this block
    %     \end{footnotesize}
    % }
\end{columns}
\end{document}